\documentclass[10pt]{beamer}
\setbeamertemplate{navigation symbols}{\insertframenumber}

\usepackage[utf8]{inputenc}
\usepackage[T1]{fontenc}%paquet de langue français%
\usepackage[french]{babel}%encodage de la police%

\usepackage{graphicx} %'affichage des images
\usepackage{enotez}

\usepackage{verbatim}
\usepackage{amsmath}
\usepackage{array,multirow,makecell}

\usepackage{listings}
\usepackage{xcolor}

\usepackage{hyperref}

\definecolor{codegreen}{named}{green}
\definecolor{codered}{named}{red}
\definecolor{codeblue}{named}{blue}
\definecolor{codeblack}{named}{black}
\definecolor{backcolour}{rgb}{0.95,0.95,0.92}

\lstdefinestyle{mystyle}{
    backgroundcolor=\color{backcolour},   
    commentstyle=\color{codegreen},
    keywordstyle=\color{codeblue},
    numberstyle=\tiny\color{codeblack},
    stringstyle=\color{codered},
    basicstyle=\ttfamily\footnotesize,
    breakatwhitespace=false,         
    breaklines=true,                 
    captionpos=b,                    
    keepspaces=true,                 
    numbers=left,                    
    numbersep=2pt,                  
    showspaces=false,                
    showstringspaces=false,
    showtabs=false,                  
    tabsize=1
}

\lstset{style=mystyle}
\usepackage[linesnumbered,ruled,french,onelanguage]{algorithm2e}

\makeatletter
\g@addto@macro{\@algocf@init}{\SetKwInput{KwOut}{Sortie}}
\makeatother


\author{Thomas Matthieu, Tellier Basile}
\date{\today}
\title{Demonstrateur : Arbre Couvrant de poids Minimal}
\usetheme{Warsaw}

\setlength{\algomargin}{1em}
\begin{document}

\begin{frame}
\maketitle

\end{frame}

\begin{frame}
\frametitle{Table des matières}
\tableofcontents
\end{frame}


\begin{frame}
\section{Notions}
\subsection{Rappel}
Rappel sur les propriétées utilisées :
  \begin{itemize}
     \item<1-> Sommets
     \item<1-> Sommets adjacents
     \item<2-> Arêtes
     \item<2-> Arêtes adjacentes
     \item<3-> Graphe
     \item<4-> Arbre
  \end{itemize}

\end{frame}

\begin{frame}
\subsection{MST}
Un arbre couvrant minimal à pour propriétées :
\newline
D'être un graphe :
  \begin{itemize}
     \item<1-> Non-Orienté
     \item<2-> Connexe
     \item<3-> Dont les arrêtes sont pondéres
  \end{itemize}

\end{frame}

\begin{frame}
\section{Applications}
Applications

\begin{examples}
  Réseau électrique
\end{examples}

\begin{examples}
  GPS
\end{examples}

\end{frame}

\begin{frame}
\frametitle{Algorithmes}
\section{Algorithmes}
\subsection{Algorithmes}

\begin{block}{Générateur aléatoire}
  Générateur de graphes aléatoires non-orienté, pondéré et connexe.
\end{block}

\begin{alertblock}{Algorithmes gloutons}
Algorithme qui réalise étape par étape,un choix optimum local, afin d'obtenir un résultat optimum global. 
\end{alertblock}

\begin{block}{Union-Find}
  Opération : Find / Union
\end{block}

\begin{block}{File de priorité}
  Enfiler / Défiler l'élément à la plus grande valeur / Tester si la file est vide
\end{block}


\end{frame}

\begin{frame}
\section{Démonstration}
\textbf{Démonstration}

\end{frame}

\begin{frame}
\section{Comparaison}
\textbf{Différence entre Prim et Kruskal}
D'être un graphe :
  \begin{itemize}
     \item<1-> Construction de l'arbre couvrant
     \item<2-> Parcours de noeud
     \item<3-> Compléxité en temps
     \item<4-> Elements du graphes
     \item<5-> Types de graphes
  \end{itemize}
\end{frame}

\begin{frame}
\section{Sources}
  Comparaison Prim Kruskal : \url{https://byjus.com/gate/difference-between-prims-and-kruskal-algorithum/}
  \newline
  Arbre Couvrant Minimal :\url{https://fr.wikipedia.org/wiki/Arbre_couvrant_de_poids_minimal}
    \newline
  Algorithme glouton : \url{https://fr.wikipedia.org/wiki/Algorithme_glouton}
    \newline
  Algorithme de Prim : \url{https://fr.wikipedia.org/wiki/Algorithme_de_Prim}
    \newline
  Algorithme de Kruskal : \url{https://fr.wikipedia.org/wiki/Algorithme_de_Kruskal}
  
\end{frame}

\end{document}
